\documentclass[11pt,a4paper]{article}
\usepackage[T1]{fontenc}
\usepackage[utf8]{inputenc}
\usepackage{lmodern}
\usepackage[margin=2.5cm]{geometry}
\usepackage{amsmath,amssymb,amsthm,mathtools}
\usepackage{booktabs}
\usepackage{microtype}
\usepackage[hidelinks]{hyperref}
\hypersetup{pdftitle={Pearson Bounds for Open-World Continual Classification}}
\setlength{\emergencystretch}{2em}
\newtheorem{theorem}{Theorem}[section]
\newtheorem{lemma}[theorem]{Lemma}
\newtheorem{proposition}[theorem]{Proposition}
\newtheorem{corollary}[theorem]{Corollary}
\theoremstyle{definition}
\newtheorem{assumption}[theorem]{Assumption}
\newtheorem{example}[theorem]{Example}
\theoremstyle{remark}
\newtheorem{remark}[theorem]{Remark}
\newcommand{\HH}{\mathcal H_k}
\newcommand{\E}{\mathbb E}
\newcommand{\Prob}{\mathbb P}
\newcommand{\norm}[1]{\left\lVert#1\right\rVert}
\newcommand{\inner}[2]{\left\langle#1,#2\right\rangle}
\newcommand{\op}{\mathrm{op}}
\newcommand{\HS}{\mathrm{HS}}
\newcommand{\D}{D_\lambda(Q;P)}
\newcommand{\Dh}{\widehat D_\lambda}
\newcommand{\pos}[1]{\left[#1\right]_+}
\DeclareMathOperator{\tr}{tr}

\title{Pearson Bounds for\\Open-World Continual Classification}
\author{}
\date{}
\begin{document}
\maketitle

\section{Classification model and regularized Pearson score}\label{sec:setup}
Let $P_1,\ldots,P_M$ denote the distributions of the regimes represented
in a fixed memory bank. An incoming batch has distribution $Q$.
For each reference $P_i$, a directed empirical Pearson-type score
is used to decide whether to associate the batch with that regime.
A small score supports association; a large score supports keeping
the incoming regime separate. An incoming regime separated from every
stored reference should remain unknown rather than be merged with an
existing regime.

We first analyze one fixed reference $P=P_i$. All pairwise results use
this notation; Section~\ref{sec:bank} gives their consequence for a fixed
finite bank and one common threshold.

\begin{assumption}[Kernel and sampling]\label{ass:setup}
Let $\HH$ be a real, separable reproducing kernel Hilbert space on a
measurable space $\mathcal X$, with measurable feature map
$\phi(x)=k(\cdot,x)$. The kernel is positive semidefinite and satisfies
\begin{equation}\label{eq:kernel}
 k(x,y)\ge0,\qquad \sup_x k(x,x)\le K,\qquad 0<K<\infty.
\end{equation}
Fix distributions $P,Q$, a positive integer $n$, and $\lambda>0$.
The reference observations $X_1,\ldots,X_n$ are i.i.d. from $P$;
the incoming observations $Y_1,\ldots,Y_n$ are i.i.d. from $Q$;
and all $2n$ observations are mutually independent. The regime
distributions are fixed. The kernel and regularization are specified
independently of the batches used to evaluate the scores.
\end{assumption}

The normalized Gaussian kernel satisfies~\eqref{eq:kernel} with $K=1$.
Pointwise nonnegativity is a separate requirement from positive
semidefiniteness. Probabilities below integrate over the reference and
incoming batches jointly; they are not conditional guarantees for every
possible realized memory sample. All logarithms are natural.

\subsection{Population and empirical moments}
For a distribution $R$, define
\begin{align}
 \mu_R&=\E_{Z\sim R}\phi(Z),&
 T_R&=\E_{Z\sim R}[\phi(Z)\otimes\phi(Z)],\label{eq:pop-moments}\\
 \widehat\mu_P&=\frac1n\sum_{r=1}^n\phi(X_r),&
 \widehat\mu_Q&=\frac1n\sum_{r=1}^n\phi(Y_r),\label{eq:emp-means}\\
 \widehat T_P&=\frac1n\sum_{r=1}^n\phi(X_r)\otimes\phi(X_r),&
 \Delta&=\mu_Q-\mu_P,\quad
 \widehat\Delta=\widehat\mu_Q-\widehat\mu_P.\label{eq:emp-moments}
\end{align}
Here $(u\otimes v)w=\inner{v}{w}_{\HH}u$.
The operator $T_P$ is an uncentered second moment; subtracting
$\mu_P\otimes\mu_P$ would define a different score.
The bounded measurable feature map ensures that these expectations
exist, as developed in Appendix~\ref{app:foundations}.
Write
\begin{equation}\label{eq:A}
 A=T_P+\lambda I,\qquad \widehat A=\widehat T_P+\lambda I.
\end{equation}
For every $f\in\HH$,
\begin{equation}\label{eq:quadratic}
 0\le\inner{f}{T_Pf}=\E_P f(X)^2\le K\norm{f}_{\HH}^2.
\end{equation}
Consequently $0\preceq T_P,\widehat T_P\preceq KI$ and
$A,\widehat A\succeq\lambda I$; both regularized operators have
bounded inverses.

\subsection{Score and association rule}
Define the population and empirical regularized Pearson-type scores
\begin{align}
 \D&=\frac12\inner{\Delta}{A^{-1}\Delta}_{\HH},\label{eq:D}\\
 \Dh&=\frac12\inner{\widehat\Delta}
                    {\widehat A^{-1}\widehat\Delta}_{\HH}.\label{eq:Dhat}
\end{align}
The semicolon indicates that $P$ supplies the second-moment operator.
The score is directed: interchanging $P$ and $Q$ generally changes it.
For a positive operator $B$ with bounded inverse, write
$\norm{v}_B^2=\inner{v}{Bv}_{\HH}$.
Appendix~\ref{app:variational} derives this regularized RKHS score
from its variational definition.

For a static threshold $\eta\ge0$, the pairwise decision is
\begin{equation}\label{eq:decision}
 \text{associate with }P\text{ if }\Dh\le\eta;
 \qquad\text{keep separate from }P\text{ if }\Dh>\eta.
\end{equation}
Static means that the threshold is specified independently of the
realized scores. It may depend on $n,K,\lambda$, confidence levels,
and an assumed population separation.
The two scenarios are
\begin{align}
 H_0:&\quad Q=P,\label{eq:H0}\\
 H_1(\delta):&\quad \D\ge\delta>0.\label{eq:H1}
\end{align}
The second scenario formalizes a non-merging regime relative to the
reference. Distinct distributions alone do not provide a specified
positive uniform margin; Appendix~\ref{app:separation} explains the
role of this assumption.

\section{Concentration of the empirical moments}\label{sec:concentration}

For $\gamma\in(0,1)$ define
\begin{equation}\label{eq:radii}
 b_n(\gamma)=\sqrt{\frac{2K}{n}}
        \left(1+\sqrt{\log\frac1\gamma}\right),\qquad
 e_n(\gamma)=\frac K{\sqrt n}
        \left(1+\sqrt{\log\frac1\gamma}\right).
\end{equation}
These deterministic confidence radii are obtained using McDiarmid's
inequality. The first concerns mean embeddings, the second an operator.

\begin{lemma}[Mean-difference concentration]\label{lem:mean}
Under Assumption~\ref{ass:setup},
\begin{equation}\label{eq:mean-event}
 \Prob\{\norm{\widehat\Delta-\Delta}_{\HH}
             \le b_n(\gamma)\}\ge1-\gamma.
\end{equation}
\end{lemma}
\begin{proof}
Set $U_r=\phi(Y_r)-\mu_Q$ and $V_r=\phi(X_r)-\mu_P$.
They are mutually independent and centered. ``Centered'' is an identity:
$\E[\phi(X)-\mu_P]=0$. It does not assume that $\mu_P=0$.
For a single distribution $R$, expansion of the squared norm gives
\begin{equation}\label{eq:feature-variance}
 \E_R\norm{\phi(Z)-\mu_R}_{\HH}^2
 =\E_R k(Z,Z)-\norm{\mu_R}_{\HH}^2\le K.
\end{equation}
Independence removes cross terms between distinct centered observations,
so
\begin{align}
 \E\norm{\widehat\Delta-\Delta}_{\HH}^2
 &=\frac1n\left(\E_Q k(Y,Y)-\norm{\mu_Q}_{\HH}^2\right)
  +\frac1n\left(\E_P k(X,X)-\norm{\mu_P}_{\HH}^2\right)\notag\\
 &\le\frac{2K}{n}.\label{eq:delta-second-moment}
\end{align}
Thus, for $F=\norm{\widehat\Delta-\Delta}_{\HH}$,
$\E F\le\sqrt{\E F^2}\le\sqrt{2K/n}$.

Replacing one observation $X_r$ by $X_r'$ changes $\widehat\Delta$ by
$-(\phi(X_r)-\phi(X_r'))/n$. The population $\Delta$ stays fixed.
The reverse triangle inequality therefore gives
\begin{align}
 |F-F'|&\le\frac1n\norm{\phi(X_r)-\phi(X_r')}_{\HH},\notag\\
 \norm{\phi(x)-\phi(y)}_{\HH}^2
 &=k(x,x)+k(y,y)-2k(x,y)\le2K.\label{eq:feature-diameter}
\end{align}
The same bound holds when replacing a $Y_r$. There are $2n$ observations,
each with sensitivity $\sqrt{2K}/n$, and hence
$\sum_{j=1}^{2n}c_j^2=4K/n$.
McDiarmid's inequality (Appendix~\ref{app:mcdiarmid}) yields
\begin{equation}\label{eq:mean-tail}
 \Prob\{F-\E F\ge t\}
 \le\exp\left(-\frac{nt^2}{2K}\right),\qquad t>0.
\end{equation}
Taking $t=\sqrt{(2K/n)\log(1/\gamma)}$ proves the claim.
\end{proof}

\begin{remark}[Expected error versus an individual deviation]
Equation~\eqref{eq:feature-variance} bounds an expectation. It does not
assert $\norm{\phi(X)-\mu_P}^2\le K$ for every $X$.
The coarser bound $\norm{\phi(X)-\mu_P}\le2\sqrt K$ is valid pointwise
and can be used to control individual summands. These two statements
serve different purposes and are compatible.
\end{remark}

\begin{lemma}[Operator concentration]\label{lem:operator}
Under Assumption~\ref{ass:setup},
\begin{equation}\label{eq:op-event}
 \Prob\{\norm{\widehat T_P-T_P}_{\op}
            \le e_n(\gamma)\}\ge1-\gamma.
\end{equation}
\end{lemma}
\begin{proof}
The Hilbert--Schmidt operators form a Hilbert space.
For $V(z)=\phi(z)\otimes\phi(z)$, the rank-one identities proved in
Appendix~\ref{app:rankone} give
\begin{equation}\label{eq:Vnorm}
 \norm{V(z)}_{\HS}=k(z,z)\le K,\qquad
 \inner{V(x)}{V(y)}_{\HS}=k(x,y)^2.
\end{equation}
Apply the Hilbert-space variance identity to their empirical mean:
\begin{equation}\label{eq:operator-variance}
 \E\norm{\widehat T_P-T_P}_{\HS}^2
 =\frac{\E_P k(X,X)^2-\norm{T_P}_{\HS}^2}{n}
 \le\frac{K^2}{n}.
\end{equation}
Thus $G=\norm{\widehat T_P-T_P}_{\HS}$ satisfies $\E G\le K/\sqrt n$.
For one replacement,
\begin{align}
 |G-G'|&\le\frac1n\norm{V(x)-V(x')}_{\HS},\notag\\
 \norm{V(x)-V(y)}_{\HS}^2
 &=k(x,x)^2+k(y,y)^2-2k(x,y)^2\le2K^2.
 \label{eq:Vdiameter}
\end{align}
There are $n$ sensitivities $\sqrt2K/n$, so their squared sum is
$2K^2/n$. McDiarmid gives
\begin{equation}\label{eq:op-tail}
 \Prob\{G-\E G\ge t\}\le\exp(-nt^2/K^2),\qquad t>0.
\end{equation}
Set $t=K\sqrt{\log(1/\gamma)/n}$ and use
$\norm{B}_{\op}\le\norm{B}_{\HS}$.
\end{proof}

\section{Null hypothesis: association with the same regime}\label{sec:null}
\begin{theorem}[Null upper bound]\label{thm:null}
Under Assumption~\ref{ass:setup}, if $Q=P$, then for every
$\alpha_0\in(0,1)$,
\begin{equation}\label{eq:null}
 \Prob\{0\le\Dh\le U_n(\alpha_0)\}\ge1-\alpha_0,
\end{equation}
where
\begin{equation}\label{eq:U}
 \boxed{U_n(\alpha_0)=\frac{b_n(\alpha_0)^2}{2\lambda}
 =\frac{K}{n\lambda}
       \left(1+\sqrt{\log\frac1{\alpha_0}}\right)^2.}
\end{equation}
Any static threshold $\eta\ge U_n(\alpha_0)$ therefore satisfies
\begin{equation}\label{eq:false-separation}
 \Prob_{Q=P}\{\Dh>\eta\}\le\alpha_0.
\end{equation}
\end{theorem}
\begin{proof}
For every realization, $\widehat A\succeq\lambda I$ implies
$0\preceq\widehat A^{-1}\preceq\lambda^{-1}I$. Hence
\begin{equation}\label{eq:null-domination}
 0\le\Dh\le\frac1{2\lambda}\norm{\widehat\Delta}_{\HH}^2.
\end{equation}
Under $Q=P$ the population mean difference is $\Delta=0$.
Lemma~\ref{lem:mean}, applied at level $\alpha_0$, then gives the
upper bound with probability at least $1-\alpha_0$.
The event $\{\Dh>\eta\}$ is contained in $\{\Dh>U_n(\alpha_0)\}$
when $\eta\ge U_n(\alpha_0)$, proving the error statement.
\end{proof}
The error in~\eqref{eq:false-separation} is failure to associate the
incoming batch with its own regime. No operator concentration is
needed under the null. The operator inequality is pointwise, so the
dependence between $\widehat A$ and $\widehat\Delta$ is allowed.

\section{Non-merging regimes: separation from a reference}\label{sec:alternative}
\begin{assumption}[Population separation]\label{ass:separation}
For the fixed ordered pair and fixed regularization, $\D\ge\delta>0$.
\end{assumption}

Both $\widehat T_P$ and $\widehat\Delta$ contain the reference data,
so they are generally dependent. For their joint use, put
$b=b_n(\alpha_1/2)$ and $e=e_n(\alpha_1/2)$ and define
\begin{equation}\label{eq:joint-event}
 \mathcal E_{b,e}=
 \{\norm{\widehat\Delta-\Delta}_{\HH}\le b,
   \ \norm{\widehat T_P-T_P}_{\op}\le e\}.
\end{equation}
The two concentration lemmas and a union bound give
\begin{equation}\label{eq:joint-probability}
 \Prob(\mathcal E_{b,e})\ge1-\alpha_1.
\end{equation}
Independence between these two error events is not required.

\begin{lemma}[Inverse-operator comparison]\label{lem:inverse}
If $\norm{\widehat T_P-T_P}_{\op}\le e$ and $c=1+e/\lambda$, then
\begin{equation}\label{eq:inverse-order}
 \widehat A^{-1}\succeq c^{-1}A^{-1}.
\end{equation}
\end{lemma}
\begin{proof}
Self-adjointness implies $\widehat T_P-T_P\preceq eI$.
Since $A\succeq\lambda I$,
\[
 \widehat A\preceq A+eI\preceq(1+e/\lambda)A=cA.
\]
Inversion reverses the order of positive operators with bounded
inverses, proving the claim. Appendix~\ref{app:order} justifies this
step without assuming commutativity.
\end{proof}

\begin{theorem}[Non-merging lower bound]\label{thm:alternative}
Under Assumptions~\ref{ass:setup} and~\ref{ass:separation}, for every
$\alpha_1\in(0,1)$,
\begin{equation}\label{eq:alternative}
 \Prob\{\Dh\ge L_n(\delta,\alpha_1)\}\ge1-\alpha_1,
\end{equation}
where
\begin{equation}\label{eq:L}
 \boxed{L_n(\delta,\alpha_1)=
 \frac{\pos{\sqrt{2\delta}-b_n(\alpha_1/2)/\sqrt\lambda}^{2}}
 {2\left(1+e_n(\alpha_1/2)/\lambda\right)}.}
\end{equation}
Here $[x]_+=\max\{x,0\}$. Any $\eta<L_n(\delta,\alpha_1)$ satisfies
\begin{equation}\label{eq:false-merging}
 \Prob\{\Dh\le\eta\}\le\alpha_1
 \quad\text{for each fixed pair with }\D\ge\delta.
\end{equation}
\end{theorem}
\begin{proof}
Use $b,e$ from~\eqref{eq:joint-event}, put $c=1+e/\lambda$, and work
on $\mathcal E_{b,e}$. The inverse comparison gives
\begin{equation}\label{eq:weighted-comparison}
 \sqrt{2\Dh}=\norm{\widehat\Delta}_{\widehat A^{-1}}
 \ge\frac1{\sqrt c}\norm{\widehat\Delta}_{A^{-1}}.
\end{equation}
Operator order holds for every vector, so it applies to the realized
$\widehat\Delta$ even though it depends on $\widehat A$.
The reverse triangle inequality in the norm induced by $A^{-1}$ gives
\[
 \norm{\widehat\Delta}_{A^{-1}}
 \ge\pos{\norm{\Delta}_{A^{-1}}
                 -\norm{\widehat\Delta-\Delta}_{A^{-1}}}.
\]
Since $A^{-1}\preceq\lambda^{-1}I$,
\[
 \norm{\widehat\Delta-\Delta}_{A^{-1}}
 \le\frac1{\sqrt\lambda}\norm{\widehat\Delta-\Delta}_{\HH}
 \le\frac b{\sqrt\lambda}.
\]
Moreover $\norm{\Delta}_{A^{-1}}=\sqrt{2\D}\ge\sqrt{2\delta}$.
Combining these facts gives
\[
 \sqrt{2\Dh}\ge\frac1{\sqrt c}
                  \pos{\sqrt{2\delta}-b/\sqrt\lambda}.
\]
Both sides are nonnegative, so squaring proves the lower bound.
Its event has probability at least $1-\alpha_1$ by the union bound.
If $\eta<L_n$, association implies $\Dh<L_n$, giving the error bound.
\end{proof}
The error in~\eqref{eq:false-merging} is an incorrect merge with a
separated reference. The mean error reduces the usable separation by
$b/\sqrt\lambda$, and the operator error introduces the factor
$1+e/\lambda$. The bound is valid for every $e\ge0$; its lower
endpoint is positive only if $\sqrt{2\delta}>b/\sqrt\lambda$.

\section{A nonempty interval for a static classification threshold}\label{sec:threshold}

Associate with the reference when $\Dh\le\eta$, and keep it separate when $\Dh>\eta$.
Under $Q=P$, keeping the batch separate is an error. Under a non-merging
regime, association with the reference is an incorrect merge. These are different probability
models, so their error levels are not added to calibrate the two tests.

\begin{theorem}[Admissible static-threshold interval]\label{cor:threshold}
For given $\alpha_0,\alpha_1\in(0,1)$, any deterministic threshold with
\begin{equation}\label{eq:threshold-interval}
 U_n(\alpha_0)\le\eta<L_n(\delta,\alpha_1)
\end{equation}
satisfies
\begin{equation}\label{eq:errors}
 \Prob_{Q=P}\{\Dh>\eta\}\le\alpha_0,\qquad
 \Prob_{\D\ge\delta}\{\Dh\le\eta\}\le\alpha_1.
\end{equation}
The second notation means each fixed pair satisfying the separation
assumption; it is not a probability distribution on pairs.
\end{theorem}
\begin{proof}
The first error event is contained in $\{\Dh>U_n(\alpha_0)\}$.
Because $\eta<L_n$, the second is contained in
$\{\Dh<L_n(\delta,\alpha_1)\}$. Apply the two theorems.
The strict inequality at the right endpoint accounts for the acceptance
rule at equality; it cannot in general be replaced by $\eta\le L_n$.
\end{proof}

\begin{proposition}[Nonempty threshold interval]\label{prop:nonempty}
Put
\[
 b_0=b_n(\alpha_0),\quad b_1=b_n(\alpha_1/2),\quad
 c=1+e_n(\alpha_1/2)/\lambda.
\]
Since $b_0>0$, the interval~\eqref{eq:threshold-interval} is nonempty
if and only if
\begin{equation}\label{eq:separation-condition}
 \boxed{\delta>\frac{(b_1+\sqrt c\,b_0)^2}{2\lambda}.}
\end{equation}
\end{proposition}
\begin{proof}
The inequality $L_n>U_n$ is equivalent to
$\pos{\sqrt{2\delta}-b_1/\sqrt\lambda}>\sqrt c\,b_0/\sqrt\lambda$.
The right side is positive, so the positive part can be removed and
the resulting positive inequality can be squared.
\end{proof}

\subsection{Common tolerance and explicit calibration}
For $\alpha_0=\alpha_1=\alpha$, choose
\begin{equation}\label{eq:eta}
 \boxed{\eta=\frac{K}{n\lambda}
       \left(1+\sqrt{\log\frac1\alpha}\right)^2.}
\end{equation}
It controls incorrect separation under the null at level $\alpha$
without a population separation assumption. The sufficient separation
for the non-merging guarantee is
\begin{equation}\label{eq:explicit-delta}
\begin{split}
 \delta>\frac{K}{n\lambda}\Bigg[&1+\sqrt{\log\frac2\alpha}\\
 &+\left(1+\sqrt{\log\frac1\alpha}\right)
 \sqrt{1+\frac{K}{\lambda\sqrt n}
                 \left(1+\sqrt{\log\frac2\alpha}\right)}
 \Bigg]^2.
\end{split}
\end{equation}
This is equivalent to nonemptiness of the displayed interval. It is not
necessary for actual successful classification:
the concentration radii may overestimate the realized errors.

\subsection{Sample size and regularization}
For fixed $K,\lambda,\alpha$ and fixed separated pair,
$b_n,e_n\to0$, $U_n\to0$, and $L_n(\delta,\alpha)\to\delta$.
Thus the certified interval eventually becomes nonempty.
For fixed confidence, the two quantities controlling usefulness are
\[
 \frac{b_n}{\sqrt\lambda}=O((n\lambda)^{-1/2}),\qquad
 \frac{e_n}{\lambda}=O((\lambda\sqrt n)^{-1}).
\]
These expressions explain why very small regularization makes the
uniform bounds demanding. If $\lambda=\lambda_n$ changes, the
population score changes, so the validity and size of an assumed
separation must also be checked. One cannot optimize the explicit
threshold while treating separation as a fixed unrelated constant.

\subsection{An explicit sufficient batch size}
\begin{corollary}[Batch size ensuring a nonempty interval]\label{cor:sample-size}
For common tolerance $\alpha$, put
\[
 a_0=1+\sqrt{\log(1/\alpha)},\qquad
 a_1=1+\sqrt{\log(2/\alpha)}.
\]
The interval is nonempty if the integer $n$ satisfies both
\begin{equation}\label{eq:batch-size}
 n\ge\left(\frac{Ka_1}{\lambda}\right)^2,
 \qquad n>\frac{K(a_1+\sqrt2\,a_0)^2}{\lambda\delta}.
\end{equation}
\end{corollary}
\begin{proof}
The first condition gives $e_n(\alpha/2)/\lambda\le1$, so $c\le2$.
Since $b_0=\sqrt{2K/n}\,a_0$ and $b_1=\sqrt{2K/n}\,a_1$, the
right side of~\eqref{eq:separation-condition} is at most
$K(a_1+\sqrt2\,a_0)^2/(n\lambda)$. The second condition makes this
strictly smaller than $\delta$.
\end{proof}

\subsection{One static threshold for a fixed memory bank}\label{sec:bank}
Let $M\ge1$, let $P_1,\ldots,P_M$ be distinct fixed reference regimes, and let
$\mathcal Q$ be an admissible family of incoming distributions.
Use the same $n,K,\lambda$ for all comparisons. Assume a uniform
non-merging margin $\delta_*>0$ such that
\begin{equation}\label{eq:bank-separation}
 D_\lambda(Q;P_i)\ge\delta_*
 \quad\text{for every }Q\in\mathcal Q
 \text{ and every }i\text{ with }Q\ne P_i.
\end{equation}
This includes all wrong-reference comparisons for known regimes
and all reference comparisons for novel regimes. Separation within
the stored bank alone would not ensure separation for arbitrary
unknown distributions.

\begin{corollary}[Known and unknown regime classification]\label{cor:bank}
Suppose each reference comparison satisfies Assumption~\ref{ass:setup}
and~\eqref{eq:bank-separation}. Scores may share the incoming batch;
reference batches need not be independent of one another. For
$\alpha\in(0,1)$, choose a common static threshold with
\begin{equation}\label{eq:bank-threshold}
 U_n(\alpha/M)\le\eta<L_n(\delta_*,\alpha/M).
\end{equation}
Let $\widehat{\mathcal A}(Q)=\{i:\widehat D_{\lambda,i}\le\eta\}$.
For $Q=P_j\in\mathcal Q$,
\begin{equation}\label{eq:known-classification}
 \Prob\{\widehat{\mathcal A}(Q)=\{j\}\}\ge1-\alpha.
\end{equation}
For $Q\in\mathcal Q$ different from every stored regime,
\begin{equation}\label{eq:unknown-classification}
 \Prob\{\widehat{\mathcal A}(Q)=\varnothing\}\ge1-\alpha.
\end{equation}
A common threshold exists whenever~\eqref{eq:explicit-delta} holds
with $\delta=\delta_*$ and $\alpha$ replaced by $\alpha/M$.
\end{corollary}
\begin{proof}
If $Q=P_j$, the null theorem bounds the event
$\widehat D_{\lambda,j}>\eta$ by $\alpha/M$. For every $i\ne j$,
the non-merging theorem bounds $\widehat D_{\lambda,i}\le\eta$
by $\alpha/M$. A union bound over these $M$ undesired decisions gives
failure probability at most $\alpha$; on their complement exactly
$j$ is accepted. For a novel $Q$, all $M$ reference comparisons
satisfy the non-merging hypothesis. A union bound on their association
events gives an empty accepted set with probability at least $1-\alpha$.
No independence between scores is used. Nonemptiness follows from
\eqref{eq:separation-condition}.
\end{proof}
The corresponding decision associates a batch when exactly one
reference is accepted, retains it as unknown when none are accepted,
and abstains when multiple references are accepted. The guarantee
is per incoming batch for a fixed bank under the stated sampling
model. It does not assume that an arbitrarily selected or adaptively
constructed memory sample has those same sampling properties.

\appendix
\section{RKHS and expectation foundations}\label{app:foundations}

\subsection{The reproducing property}
For $\phi(x)=k(\cdot,x)$,
\[
 f(x)=\inner{f}{\phi(x)}_{\HH},\qquad
 \inner{\phi(x)}{\phi(y)}_{\HH}=k(x,y),\qquad
 \norm{\phi(x)}_{\HH}^2=k(x,x).
\]
Cauchy--Schwarz implies $|k(x,y)|\le\sqrt{k(x,x)k(y,y)}\le K$
and $|f(x)|\le\sqrt K\norm{f}_{\HH}$.
A Hilbert-valued random element has a Bochner expectation when it is
strongly measurable and its expected norm is finite. Our assumptions
give these properties, since $\norm{\phi(X)}\le\sqrt K$ and
$\norm{\phi(X)\otimes\phi(X)}_{\HS}\le K$.
Linear functionals pass through these expectations, which justifies
identities such as $\E f(X)=\inner{f}{\mu_P}$.

\subsection{Boundedness and positivity of the second moment}
For $f\in\HH$,
\[
 0\le\inner{f}{T_Pf}=\E_P f(X)^2
 \le K\norm{f}_{\HH}^2.
\]
This proves $0\preceq T_P\preceq KI$.
The empirical version is proved by replacing the expectation with a
finite average. Consequently $A^{-1}$ and $\widehat A^{-1}$ exist,
and each has operator norm at most $1/\lambda$.

\subsection{Why the score can be zero for distinct distributions}
Since $A\preceq(K+\lambda)I$,
\begin{equation}\label{eq:score-MMD}
 \frac{\norm{\Delta}^2}{2(K+\lambda)}
 \le\D\le\frac{\norm{\Delta}^2}{2\lambda}.
\end{equation}
Thus zero score is equivalent to equal mean embeddings. A noncharacteristic
kernel can map different distributions to the same mean embedding.
Positive score separation must therefore be assumed or established
using properties of the kernel and the admissible distributions.

\section{Centering, variance, and independence}\label{app:variance}

\subsection{The Hilbert-space variance identity}
Let $W$ be a square-integrable random element of a real Hilbert space
and $\mu=\E W$. Expand:
\begin{align*}
 \E\norm{W-\mu}^2
 &=\E\norm{W}^2-2\E\inner{W}{\mu}+\norm{\mu}^2\\
 &=\E\norm{W}^2-2\inner{\E W}{\mu}+\norm{\mu}^2\\
 &=\E\norm{W}^2-\norm{\mu}^2.
\end{align*}
The middle term does not vanish: it equals $-2\norm{\mu}^2$.
It combines with the last term to leave a subtraction.
No assumption that the uncentered observation has mean zero is used.

\subsection{Independent empirical averages}
If $W_1,\ldots,W_n$ are independent copies of $W$, let
$R_r=W_r-\mu$. Independence gives, for $r\ne s$,
\[
 \E\inner{R_r}{R_s}=\inner{\E R_r}{\E R_s}=0.
\]
Hence
\[
 \E\left\lVert\frac1n\sum_rR_r\right\rVert^2
 =\frac1{n^2}\sum_{r,s}\E\inner{R_r}{R_s}
 =\frac1n\left(\E\norm{W}^2-\norm{\mu}^2\right).
\]

\subsection{An example separating a pointwise bound from an expectation}
Consider two points $a,b$ with $k(x,y)=\mathbf1\{x=y\}$ and
$P(a)=p$, $P(b)=1-p$. The feature vectors are orthonormal, $K=1$,
and
\[
 \norm{\phi(a)-\mu_P}^2=2(1-p)^2,\qquad
 \E\norm{\phi(X)-\mu_P}^2=2p(1-p)\le\frac12.
\]
At $p=0.01$, the first quantity is $1.9602>K$, but its expectation
is $0.0198<K$. An individual large error does not contradict a small
second moment. In particular, the exact second-moment estimate cannot
be substituted for a pointwise summand bound in a concentration theorem.

\subsection{What dependence would change}
If $U=\widehat\mu_Q-\mu_Q$ and $V=\widehat\mu_P-\mu_P$, then always
\[
 \E\norm{U-V}^2=\E\norm{U}^2+\E\norm{V}^2
                      -2\E\inner{U}{V}.
\]
Centering alone does not remove the last term. Independent windows
do remove it. Dependence within a window likewise changes the
variance of that window's mean. The classification results require the independence assumptions in
Assumption~\ref{ass:setup}.

\section{Rank-one operators and Hilbert--Schmidt norms}\label{app:rankone}

Let $(e_j)$ be an orthonormal basis of $\HH$. A Hilbert--Schmidt
operator $B$ has $\norm{B}_{\HS}^2=\sum_j\norm{Be_j}^2$.
For rank-one operators, Parseval's identity gives
\begin{align*}
 \norm{u\otimes v}_{\HS}^2
 &=\sum_j\norm{\inner{v}{e_j}u}^2
 =\norm{u}^2\sum_j|\inner{v}{e_j}|^2
 =\norm{u}^2\norm{v}^2,\\
 \inner{u\otimes v}{a\otimes b}_{\HS}
 &=\inner{u}{a}\inner{v}{b}.
\end{align*}
Taking $u=v=\phi(x)$ and $a=b=\phi(y)$ gives~\eqref{eq:Vnorm}.
Expanding the squared difference then gives~\eqref{eq:Vdiameter}.
For any unit vector $f$, extending $f$ to an orthonormal basis shows
$\norm{Bf}^2\le\norm{B}_{\HS}^2$; taking a supremum proves
$\norm{B}_{\op}\le\norm{B}_{\HS}$.

It is the error $\widehat T_P-T_P$ that is Hilbert--Schmidt. In an
infinite-dimensional RKHS, $I$ need not be Hilbert--Schmidt, and we
do not claim that $A=T_P+\lambda I$ is Hilbert--Schmidt. We use
operator norms and order for the regularized operators instead.

\section{McDiarmid's inequality and constant bookkeeping}\label{app:mcdiarmid}

Let $Z_1,\ldots,Z_N$ be independent, and let $f$ be a measurable
real function such that changing coordinate $j$ changes $f$ by at
most a deterministic $c_j$. Then, for $t>0$ and $\sum_j c_j^2>0$,
\begin{equation}\label{eq:mcdiarmid}
 \Prob\{f-\E f\ge t\}
 \le\exp\left(-\frac{2t^2}{\sum_{j=1}^Nc_j^2}\right).
\end{equation}
It follows that, with probability at least $1-\gamma$,
\[
 f\le\E f+\sqrt{\frac12\sum_jc_j^2\log(1/\gamma)}.
\]
The squared sum must count \emph{all} independent input coordinates.
For the mean difference there are $2n$; for the reference operator
there are $n$.
\begin{center}
\begin{tabular}{llll}
\toprule
Function & Inputs & Sensitivity & Squared sum\\
\midrule
$\norm{\widehat\Delta-\Delta}$ & $2n$ & $\sqrt{2K}/n$ & $4K/n$\\
$\norm{\widehat T_P-T_P}_{\HS}$ & $n$ & $\sqrt2K/n$ & $2K^2/n$\\
\bottomrule
\end{tabular}
\end{center}
For a two-sided scalar deviation around the expectation, a union bound
over the upper and lower tails introduces a factor $2$ in probability,
or $\log(2/\gamma)$ in the radius. Our main lemmas need only an
upper tail for an already nonnegative error norm, so they use
$\log(1/\gamma)$ instead.

\section{Operator order and the weighted norm}
\label{app:order}

\subsection{Order and inversion}
For self-adjoint operators, $B\preceq C$ means
$\inner{v}{Bv}\le\inner{v}{Cv}$ for every $v$.
If $0\prec B\preceq C$ with bounded inverses, let
$S=C^{-1/2}BC^{-1/2}$. Then $0\prec S\preceq I$, so its spectral
values are in $(0,1]$ and $S^{-1}\succeq I$. Consequently,
\[
 B^{-1}=C^{-1/2}S^{-1}C^{-1/2}\succeq C^{-1}.
\]
This proves order reversal without multiplying noncommuting inequalities
in an invalid order. Apply this to $B=\widehat A$ and $C=cA$ in Lemma~\ref{lem:inverse}.

\subsection{The reverse triangle inequality in the weighted norm}
Since $A^{-1}$ is positive with bounded inverse,
\[
 \norm{v}_{A^{-1}}=\norm{A^{-1/2}v}_{\HH}
\]
is a norm. Its triangle inequality is the Hilbert-space triangle
inequality after applying the linear map $A^{-1/2}$. Therefore
\[
 \norm{u}_{A^{-1}}\ge
 \norm{v}_{A^{-1}}-\norm{u-v}_{A^{-1}}.
\]
The left side is nonnegative, so the positive part can be taken on
the right. This is the positive part in the non-merging bound.
Also $A^{-1}\preceq\lambda^{-1}I$ implies
$\norm{u-v}_{A^{-1}}\le\norm{u-v}_{\HH}/\sqrt\lambda$.

\subsection{Why shared reference observations are allowed}
For any two events, a union bound gives
$\Prob(E_1\cap E_2)\ge1-\Prob(E_1^c)-\Prob(E_2^c)$.
Once both error bounds hold, operator order is an inequality for every
vector on that outcome. It can be applied to the realized
$\widehat\Delta$ without assuming that it is independent of
$\widehat A$. In particular, the argument does not factor an
expectation involving these estimators.

\subsection{An explicit example where the estimators are coupled}
Take $\mathcal X=\{0,1\}$ and the scalar kernel $k(x,y)=xy$.
It is positive semidefinite and nonnegative, with $K=1$; identify
$\phi(x)=x$ in a one-dimensional feature space. Let $n=1$,
$P=Q=\operatorname{Bernoulli}(0.2)$, $X,Y$ independent, and
$\lambda=1$. Then
\[
 \widehat T_P=X,\quad \widehat\Delta=Y-X,\quad
 \widehat A^{-1}=\frac1{1+X},\quad
 \Dh=\frac{(Y-X)^2}{2(1+X)}.
\]
The events $(X,Y)=(0,1)$ and $(1,0)$ each have probability $0.16$.
Their scores are $1/2$ and $1/4$, respectively; the other scores vanish.
Therefore
\[
 \E\Dh=0.16\left(\frac12+\frac14\right)=0.12,
 \qquad
 \frac12\E[\widehat\Delta^2]\E[\widehat A^{-1}]
 =\frac12(0.32)(0.9)=0.144.
\]
The failed factorization demonstrates actual dependence. Our null
domination and joint-event operator argument remain valid because
neither uses this factorization.

\section{Separation, identifiability, and sufficient conditions}
\label{app:separation}

For a fixed finite collection with injective mean embeddings and fixed
$\lambda$, every distinct ordered pair has positive score. The minimum
over finitely many such positive scores is positive. This supports a
finite-bank separation assumption.
It does not imply a positive minimum over all distinct distributions.
To see why, fix $P,R$ and let $Q_\varepsilon=(1-\varepsilon)P+\varepsilon R$.
Then
\[
 \mu_{Q_\varepsilon}-\mu_P=\varepsilon(\mu_R-\mu_P),\qquad
 D_\lambda(Q_\varepsilon;P)=\varepsilon^2D_\lambda(R;P)\longrightarrow0.
\]
This example applies even to characteristic Gaussian kernels. A uniform
positive separation for unknown regimes therefore restricts the allowed
alternatives; it is not a consequence of $P\ne Q$.

The static-threshold separation criterion~\eqref{eq:separation-condition} is necessary
and sufficient for nonemptiness of the interval defined by our chosen
$U_n,L_n$. It is only a sufficient condition for small actual error.
Failure of this criterion means the displayed bounds do not certify
the desired threshold, not that classification is impossible.

\section{Variational derivation of the classification score}\label{app:variational}
For a reference distribution $P$, the regularized RKHS Pearson
functional is
\begin{equation}\label{eq:variational}
 \sup_{f\in\HH}\left\{
 \E_Qf-\E_Pf-\frac12\E_P f^2-\frac\lambda2\norm{f}_{\HH}^2
 \right\}.
\end{equation}
The reproducing property identifies the expression in braces with
$\inner{f}{\Delta}-\frac12\inner{f}{Af}$. Completing the square gives
\[
 \inner{f}{\Delta}-\frac12\inner{f}{Af}
 =\frac12\inner{\Delta}{A^{-1}\Delta}
  -\frac12\norm{f-A^{-1}\Delta}_{A}^{2}.
\]
The final squared norm is nonnegative and vanishes only at
$f=A^{-1}\Delta$. The supremum is therefore exactly~\eqref{eq:D}.
Replacing population expectations by empirical averages yields
\eqref{eq:Dhat}, with maximizer $\widehat A^{-1}\widehat\Delta$.
The factor $1/2$ is part of the normalization, and the restriction to
$\HH$ and the regularization are part of the score itself. The
separation assumption refers to this score, not to an unrestricted
density-ratio functional.

\section{Expansion of the static-threshold condition}\label{app:threshold}
At common tolerance $\alpha$, write
\[
 a_0=1+\sqrt{\log(1/\alpha)},\qquad
 a_1=1+\sqrt{\log(2/\alpha)}.
\]
Then
\[
 b_0=\sqrt{\frac{2K}{n}}a_0,\qquad
 b_1=\sqrt{\frac{2K}{n}}a_1,\qquad
 c=1+\frac{K}{\lambda\sqrt n}a_1.
\]
Substituting in~\eqref{eq:separation-condition},
\begin{align*}
 \frac{(b_1+\sqrt c\,b_0)^2}{2\lambda}
 &=\frac1{2\lambda}\frac{2K}{n}
   \left(a_1+\sqrt{1+\frac{K}{\lambda\sqrt n}a_1}\,a_0\right)^2\\
 &=\frac{K}{n\lambda}
   \left(a_1+a_0\sqrt{1+\frac{K}{\lambda\sqrt n}a_1}\right)^2.
\end{align*}
Expanding $a_0,a_1$ gives~\eqref{eq:explicit-delta}. The strict
inequality means that the lower confidence bound for a non-merging
regime lies strictly above the null upper confidence bound. These
are bounds for two different population scenarios. Their separation
certifies that one predetermined threshold can implement both
classification decisions at the stated error levels.

\end{document}
